\documentclass[12pt]{article}
\usepackage[margin=1in]{geometry}
\usepackage{enumitem}
\usepackage{titlesec}
\usepackage{parskip}
\usepackage{amsmath}
\usepackage{graphicx}
\usepackage{hyperref}
\usepackage{fontenc}

\title{\textbf{Putnam 1993 RG Notes}\\
\large Putnam, Robert D. \textit{Making Democracy Work: Civic Traditions in Modern Italy}.\\
Princeton: Princeton University Press, 1993.}
\author{}
\date{}

\begin{document}

\maketitle

\section{Main Idea}



%--------------------------------------------------
\section{General Notes}
%--------------------------------------------------

\begin{itemize}[leftmargin=*, itemsep=4pt]
    \item Central question: ``What are the conditions for creating strong, responsive, effective representative institutions?'' Comparison of regional Italian bureaucracies as an extended case, exploiting the simultaneous creation of fifteen new regional governments in the 1970s: homogenous constitutional authority, but drastically different policy results. Argument is that differing social and political contexts lead to these results
    \item Transformation of existing metrics of institutional performance: we do not merely need the government to make decisions, but we further need them to \emph{act} on those decsions. Agreement itself (a social choice theoretic outcome) is not sufficient to demonstrate institutional success. Model is:
    
    \begin{align*}
        \text{Societal demands} & \implies \text{Political interaction} \\
        & \implies \text{Government} \\
        & \implies \text{Policy choice} \\
        & \implies \text{Implementation}
    \end{align*}
    ``A high-performance democratic institution must be both responsive and effective: sensitive to the demands of its constituents and effective in using limited resources to address those demands.''
    \item Core transition in institutional authority came with the ability to establish their own patronage system: here called ``safeguarding public morals,'' the new institutions (regional) were ultimately allowed to create their own licensing, permitting, and regulatory efforts that could exploit (and be exploited by) local actors.
    \item With regionalization came a new generation of regional elites, defined by a ``passing of the torch'' to those with lesser name recognition and hereditary status and instead self-made wealth and status
    \item Any metric of government performance should satisfy the following:
    \begin{enumerate}
        \item \underline{Comprehensive}: must consider the wide range of potential roles of government, and act accordingly
        \item \underline{Internally consistent}: must consider the fact that governments can be measured along different frameworks and dimensions, and must approximately hold that rank
        \item \underline{Reliable}: should be generalizable rather than highly volatile: can/should change as performance does, but should not be super-overly-reactive in the absence of meaningful reasons to redirect
        \item \underline{Correspond to the objectives and evaluations of the institution's protagonists and constituents}: must actually reflect the needs/desires/metrics of the people subjected to them
    \end{enumerate}
    \item One driving theory of differential efficacy stems from the civic context: the ``civic virtue'' of the people, as historical normative theory would call it. This was often in contrast with republicanism, but had re-swept Western thought in the mid-20th century. 11As the proportion of nonvirtuous citizens increases significantly, the ability of liberal societies to function successfully will progressively diminish."
    \item Some possible dimensions of civic capacity:
    \begin{itemize}
        \item Civic engagement: are citizens active in the political process? Not necessarily a 1:1 here, as action can by non-virtuous or otherwise fueled by solely self-interest. Does not defeat a collective action problem -- most citizens still look for short-term maximization of the immediate family
        \item Political equality: do citizens interact as equals, with themselves and with their government officials?
        \item Solidatiry, trust, and tolerance: what do the bonds between citizens look like? Are they strong even across heterogenous groups, identities, and backgrounds?
        \item Associations: are there distinctive social structures and norms? Are there civil and political organizations, particularly those founded by citizens rather than imposed by the government? (This is a bit of an Ostromian thought as well, regarding endogenous institutions) Also civic institutions that are entirely apolitical -- think later Putnam ``Bowling Alone''
    \end{itemize}
    \item A significant factor is social trust: citizens in civic communities trust each other and expect fair dealings, and similarly trust the government's ability to achieve high standards in exchange for compliance (a heightened social contract). ``Collective life in the civic regions is eased by the expectation that others will probably follow the rules. Knowing that others will, \emph{you} are more likely to go along, too, thus fulfilling \emph{their} expectations.''
    \item In the absence of such systems of trust and reciprocal enforcement of collective action problems, less civic communities turn towards forces of coercion (i.e., the police). This is the only alternative solution to the Leviathan. ``In the absence of solidarity and self-discipline, hierarchy and force provide the only alternative to anarchy.'' This also results in a paradox whereby more individualistic societies are actually more state-dependent for enforcing order.
    \item Heightened social trust allowed for financial growth and significant expansion of trade; in turn, the prosperity these returned built even further social trust (and trust in the civic institutions). A sign/portion of this was the rise of credit and loanmaking.
    \item Positive relationship between civic traditions and civic community; directionality is a bit hard to establish, but plausible that having an established set of traditions allows for social coherence which in turn lays the ground for civic community.
    \item \emph{``Economics does not predict civics, but civics does predict economics, better indeed than economics itself.''}
    \item Social capital -- and ``rotating credit associations'' -- can be a solution to the collective action problem, by means of building trust in one another and subsequently allowing for higher-risk activities and investments. Rotating credit association members buy into a fund which is distributed later to one member. Subsequently, social capital ``is the key to making democracy work.''
\end{itemize}


\section{Important Specific Factual Claims}

\begin{itemize}[leftmargin=*, itemsep=4pt]
    \item Arturo Israel: It is easier to build a road than to build an organization to maintain that road
    \item Pre-regional reformation in the 1970s, public policy in Italy was almost entirely the result of centralized oversight and public affairs, mediated by heavy patronage influences. This centralization was maintained throughout various iterations of the national Italian Government
    \item Distinction between North and South in Italy: the North was more concerned by central threats to authority, while the South was more concerned by central threats to funding. North favored ``horizontal" collective action across a wide front while South favored ``vertical'' action through private petitions, etc.
    \item ``The first test for any new political institution is that it must engage the aspirations and harness the ambitions of serious politicos.''
    \item Broad depolarization of Italian politics (at least at the local level) through and beyond this period of reformation; shift of local politics from ideological extremes to pragmatic policymaking. Much of this may be explained by institutional socialization, as the new institutional structure resulted in more moderates holding power and facilitated moderating interactions between ideologically opposed actors
    \item Ties back to early medieval times, whereby self-governance (in some form) arose out of serfdom in the north more than in the south of Italy -- ultimately became the communes, first as voluntary organizations and later as actual governance. Still had a sense of self-enforcement, where illicit activity was met with ostracization. Beginning of the oaths of mutual security; rise of the notary, which became a significant institution in organizing civic life
    \item ``Industrial districts'' with high rates of production/manufacturing etc. Characterized by a mix of vigorous competition between firms to innovate, but cooperation to establish common administrative/legal/financing/research outcomes. This creates a mix of institutional adaptiveness and economic strength that ultimately fueled a great deal of economic development in late-20th-century Italy.
    \item Social capital as a ``moral resource'': must be used in order to secure its presence and power, and actually decreases in strength if not exercised. Thus, perpetuates greater activity (when appropriately set)
\end{itemize}


\section{Connections}

\begin{itemize}
    \item Intense social and economic development in post-WW2 Italy sets a similar context a la Huntington 1968, particularly as political institutions did \textit{not} keep up with this level of development. 
    \item Directly references Huntington's concepts of questioned authority / development of autonomy away from previously-tribalized structures: regional autonomy in Italy as a prime example of that, as people began to question patronage/heritage-based systems, etc.
    \item Another reference to ``input'' vs. ``output'' of government a la Dahl 1972: government input is the processes and institutions which exist and are constructed to govern, and output are the actual laws and policies that come with them
    \item Other scholars (Edmund Burke, Bariel Almond, back to Cicero) cite political consensus as a foundation for good government; Putnam disagrees. Rather, civic communities / networks allow for dissent to be appropriately handled throughout society without driving signficant schisms. Evidence in Italy does not suggest that institutional success was correlated with political homogeneity
    
\end{itemize}


\section{My Critiques}

\begin{itemize}
    \item None. I love Putnam and he is near and dear to my heart
\end{itemize}


\end{document}
